\scanpage{177}
\begin{center}- 28 -\end{center}\hfill IX

\uline{Lemme} 3.1.3. \uline{Soient} $S$ \uline{un schéma}, $K$ \uline{un} S-\uline{schéma quasi-fini}, \uline{de présentation finie}, \uline{séparé et plat sur} $S$, $s\in S$, $t\in T$ \uline{une générisation de} $s$.
\uline{Alors on a}
\[
 \operatorname{rang}K_s\leq\operatorname{rang}K_t\qquad.
\]
\uline{Si rang} $K_\eta\leq1$ \uline{pour tout point maximal de} $S$, \uline{alors}
$K\longrightarrow S$ \uline{est une immersion ouverte} (\uline{donc un isomorphisme s'il admet une section}).

Pour la première assertion, on se ramène par localisation
stricte au cas où $S$ est strictement local. Décomposant alors $K$ en
somme d'un schéma fini $K^f$ et d'un schéma à fibre spéciale vide, on
aura
\[
 \operatorname{rang}K_s=\operatorname{rang}K_s^f=\operatorname{rang}K_t^f
 \leq\operatorname{rang}K_t\qquad,
\]
la deuxième égalité provenant du fait que $K^f$ est plat, fini et de
présentation finie sur $S$, donc localement libre sur $S$.

Pour la deuxième assertion, ce qui précède montre déjà que
pour tout $s\in S$ on a $\operatorname{rang}K_s\leq1$, donc $K_s$ est vide ou isomorphe à
$\operatorname{Spec}k(s)$. Par suite $K\longrightarrow S$ est un monomorphisme (SGA 3 $\mathrm{VI_B}$ 2.11),
et étant plat et de présentation finie c'est une immersion ouverte
(EGA IV 17.9.1), ce qui achève la démonstration.

\uline{Remarque} 3.1.4. Dans 3.1 a) et la suite de la proposition 3.1, on ne
donne aucun renseignement sur ${}_nK_s$ lorsque $n$ est une puissance de $p$
(lorsque $p>1$). Lorsque $G$ est lisse (et $S$ localement noethérien ou
$H$ de présentation finie sur $S$) tout au moins, on peut donner un tel
renseignement, par essentiellement la même méthode que dans la
\figure{\begin{tikzpicture}[x=.7cm,y=.7cm]
\draw[gray] (0,.1)--(.3,.12) (.3,.5)--(.3,.15) (.5,.08)--(.7,.1)
(1.7,.1)--(1.75,.01) (2,.02)--(2.25,.15)--(2.2,.3)
(3,.1)..controls(3.3,.18)and(3.5,.08)..(3.65,.13)
(4.3,.12)--(4.55,.14)--(4.6,.05)
(5,.08)--(5.1,.05) (5.5,.06)--(5.55,.17)
(6,.05)--(6.1,.06) (6.7,.03)--(6.75,.13)
(7.5,.02)--(7.7,.13) (7.7,.03)--(8.1,.04)
(10.3,-.03)..controls(10.6,-.1)and(10.9,-.13)..(11.1,0)
(12.2,.03)--(12.5,.2) (12.6,-.06)--(12.9,-.06)
(13.3,-.1)--(13.25,.02)--(13.35,.09) (13.8,-.22)--(13.85,.4);
\end{tikzpicture}}
\margin{346}
\margin{346}
\margin{176}
